\section{Summary}\label{sec:txid-summary}

\begin{outline}
  \item touch only very briefly, as the introduction already makes the argument: tracing in payments is often discussed as if it were a backdoor in end-to-end encrypted messaging, but it is neither technically nor socially the same problem
  \item a backdoor is a master capability (everything, any time, unnoticed); here one user, one epoch, only after that user's grant, non-targeted users never exposed
  \item a message has meaning only to sender and recipient, so any reader is an intrusion; a payment moves value in a system that already answers to regulators; the realistic alternative is the status quo, banks seeing everything in the clear
\end{outline}

\rebekah{more citations for the financial requirements (tracing obligations,
regulators) earlier in this chapter, so the summary can refer back to them}

\begin{outline}
  \item scoping came from moving the question from who may see a transaction to how narrowly a capability can be granted
  \item epoch-local index space (not a counter or per-user state) gives a permission a boundary in time, an owner, and a fixed cost independent of ledger size
  \item same shape elsewhere: tax reporting over one period, a court order for one account, a compliance check for one counterparty; the design question is the smallest capability that meets the need, and whether its use is visible to the person concerned
\end{outline}

\begin{outline}
  \item future work: grants are issued by users, so compelled disclosure relies on the legal system; a threshold auditor with the techniques of \Cref{chap:tsps} would split the authority between institutions
  \item simulator rewinds, so the result is standalone; straight-line extraction would give a composable one at the cost of larger proofs
  \item the algebraic tag proof is the largest item per transaction; the tag itself is 576 bytes
  \item a policy for choosing $\Max$ (privacy of volume against capacity) is a deployment question as much as a cryptographic one
  \item algebraic PRFs that fit Sigma protocols, like Dodis-Yampolskiy, seem to need $q$-type assumptions (\Cref{sec:txid-ft}); worth asking whether a different primitive avoids them, since one-way functions alone are not enough (their outputs need not look random)
  \item a redesign of how keys are derived: a long-term key derives epoch keys, which derive tags; only the current epoch key is released to an auditor; it must not give previous or future epoch keys, or previous or future tags, and tags must not give any key; if each stage can be one way like this and still have short proofs, the generic protocol gets a clean generic instantiation
\end{outline}
